%% ch02-pipeline.tex — The Security Pipeline
%% JA4proxy Technical Reference Manual

\chapter{The Security Pipeline}
\index{pipeline}
\label{ch:pipeline}

\begin{chapterquote}
The pipeline is the ordered sequence of checks applied to every TLS connection.
Every layer runs in microseconds. Signal collection layers run in parallel.
The pipeline produces exactly one action per connection.
\end{chapterquote}

\newthought{The pipeline has three zones}: the \textbf{fast-path bypass zone} (layers 0–2),
which short-circuits immediately on certain conditions; the \textbf{signal collection zone}
(layers 3–10), where eight checks run concurrently; and the \textbf{decision zone}
(layers 11–12), which aggregates signals into a score and maps the score to an action.

\section{Pipeline Overview}
\index{pipeline!overview}

\begin{fullwidth}
\begin{table}[h]
\caption{Complete pipeline layer reference}
\label{tab:pipeline-layers}
\begin{tabular}{@{}llllp{5.5cm}@{}}
\toprule
\textbf{Layer} & \textbf{Name} & \textbf{Zone} & \textbf{Output} & \textbf{Trigger condition} \\
\midrule
\bypassrow{0} & \bypassrow{IP trust \& normalisation} & \bypassrow{Fast path} & \bypassrow{—} & \bypassrow{Every connection} \\
\bypassrow{0b} & \bypassrow{Static IP allowlist} & \bypassrow{Fast path} & \bypassrow{\badgeallow} & \bypassrow{IP in static allowlist} \\
\bypassrow{0c} & \bypassrow{GeoIP country block} & \bypassrow{Fast path} & \bypassrow{\badgeblock} & \bypassrow{Country in blacklist} \\
\bypassrow{0d} & \bypassrow{CIDR block} & \bypassrow{Fast path} & \bypassrow{\badgeblock} & \bypassrow{IP in blocked subnet} \\
\bypassrow{0e} & \bypassrow{Spamhaus DROP/EDROP} & \bypassrow{Fast path} & \bypassrow{\badgeblock} & \bypassrow{IP in DROP/EDROP trie} \\
\bypassrow{1} & \bypassrow{ALPN browser bypass} & \bypassrow{Fast path} & \bypassrow{\badgeallow} & \bypassrow{ALPN = h2 or h1} \\
\bypassrow{1b} & \bypassrow{JA4 whitelist} & \bypassrow{Fast path} & \bypassrow{\badgeallow} & \bypassrow{JA4 in whitelist SET} \\
\bypassrow{1c} & \bypassrow{mTLS client cert} & \bypassrow{Fast path} & \bypassrow{\badgeallow} & \bypassrow{Valid client certificate} \\
\bypassrow{2} & \bypassrow{JA4 blacklist} & \bypassrow{Fast path} & \bypassrow{\badgeblock} & \bypassrow{JA4 in blacklist SET} \\
\signalrow{3} & \signalrow{TLS enforcement} & \signalrow{Signal} & \signalrow{Signal or \badgeblock} & \signalrow{TLS version/cipher check} \\
\signalrow{4} & \signalrow{SNI analysis} & \signalrow{Signal} & \signalrow{Signal} & \signalrow{Every connection} \\
\signalrow{5} & \signalrow{TCP \& connection behaviour} & \signalrow{Signal} & \signalrow{Signal} & \signalrow{Every connection} \\
\signalrow{6} & \signalrow{ASN classification} & \signalrow{Signal} & \signalrow{Signal} & \signalrow{Every connection} \\
\signalrow{7} & \signalrow{FCrDNS enrichment} & \signalrow{Signal} & \signalrow{Signal} & \signalrow{Every connection} \\
\signalrow{8} & \signalrow{Beaconing detector} & \signalrow{Signal} & \signalrow{Signal} & \signalrow{Every connection} \\
\signalrow{9} & \signalrow{AbuseIPDB score} & \signalrow{Signal} & \signalrow{Signal} & \signalrow{Every connection} \\
\signalrow{10} & \signalrow{RDAP org reputation} & \signalrow{Signal} & \signalrow{Signal} & \signalrow{Every connection} \\
11 & Attacker attribution & Analysis & Signal & Cross-IP JA4/JA4X correlation \\
12 & Behavioural analysis & Analysis & Signal & Sequential/coordinated patterns \\
13 & Risk scorer & Decision & Score 0–100 & Always \\
14 & Action decider & Decision & Action & Always \\
\bottomrule
\end{tabular}
\end{table}
\end{fullwidth}

\section{The Fast-Path Bypass System}
\index{bypass!fast path}
\index{bypass!configuring}

\subsection{How Bypasses Work}

Each bypass condition is independently configurable under the \configkey{security\_policy}
configuration section. When a bypass condition is met and its bypass is enabled, the
pipeline short-circuits immediately — the signal collection and scoring layers never
execute, and no Redis signal-collection writes are made.

All bypass conditions are toggleable. Disabling an ALLOW bypass causes those connections
to pass through the full scorer. Disabling a BLOCK bypass routes those connections through
the scorer, where the dial controls whether they are actually blocked.

\begin{warningbox}[Bypass Modification Risk]
The proxy emits a startup \code{WARN} log entry and increments a Prometheus counter for
every high-risk bypass that is disabled. Operators should monitor
\code{ja4proxy\_bypass\_total} to verify bypass layer traffic volumes are as expected.
\end{warningbox}

\subsection{Layer 0 — IP Trust and Normalisation}
\index{PROXY protocol}
\index{IP normalisation}
\index{trusted upstream CIDR}

Every connection enters layer 0. This layer:
\begin{enumerate}
  \item Verifies the connection originates from a trusted upstream CIDR (the HAProxy address range)
  \item Reads the PROXY protocol v2 header to extract the real client IP
  \item Canonicalises the IP to its compressed string representation
  \item Checks whether the IP is in the in-process ban cache (fast reject before Redis)
\end{enumerate}

If the sender is not in the trusted upstream CIDR, the connection is rejected immediately.
This prevents IP spoofing via injected PROXY protocol headers.\sidenote{Without trusted
upstream CIDR verification, an attacker who can reach port 8080 directly could inject a
PROXY header claiming any client IP.}

\subsection{Layer 0b — Static IP Allowlist}
\index{bypass!static IP allowlist}
\index{static IP allowlist}

Configuration key: \configkey{security\_policy.static\_ip\_allowlist.enabled}

If the client IP is found in the static IP allowlist (configured in
\configkey{config/proxy.yml} and hot-reloadable), the connection is immediately allowed
without scoring. Typical use: trusted monitoring systems, CI test runners, internal
health checkers.

\subsection{Layer 0c — GeoIP Country Block}
\index{GeoIP}
\index{country blocking}
\index{bypass!country blacklist}

Configuration key: \configkey{security\_policy.country\_blacklist\_bypass.enabled}

Country lookup uses the IP2Location LITE database in conjunction with a Redis-backed
country blacklist. Countries in \configkey{geoip.country\_blacklist} trigger an immediate
BLOCK. This check runs before any TLS parsing — the country decision is based on IP alone.

\begin{notebox}[Country Allowlist vs Blacklist]
\configkey{geoip.country\_whitelist\_enabled} and \configkey{geoip.country\_blacklist\_enabled}
are independent toggles. The allowlist mode (only listed countries are permitted) is
distinct from the blacklist mode (listed countries are blocked). Both can be active
simultaneously; the whitelist check runs first.
\end{notebox}

\subsection{Layer 0d — CIDR Block}
\index{CIDR blocking}

Subnet blocks are maintained in Redis and refreshed every 30 seconds into an in-process
trie. This layer is the enforcement point for RDAP block expansion (Phase~11) and
manually added subnet bans. Matching an entry produces an immediate BLOCK without
scoring.

\textbf{IPv6:} The trie handles both IPv4 and IPv6 prefixes natively. Block expansion
never adds prefixes shorter than /24 (IPv4) or /48 (IPv6).

\subsection{Layer 0e — Spamhaus DROP/EDROP}
\index{Spamhaus!DROP}
\index{Spamhaus!EDROP}
\index{bypass!Spamhaus}

Configuration key: \configkey{security\_policy.spamhaus\_bypass.enabled}

Spamhaus DROP and EDROP lists are fetched at startup and refreshed every
\configkey{spamhaus.refresh\_interval} seconds (default 3600). They are stored
in an in-process trie for O(log n) lookup.\sidenote{The Spamhaus DROP list contains
IP ranges that have been hijacked or directly allocated to cybercriminals. False
positive rate is near zero — these ranges are never used for legitimate traffic.}

When disabled, Spamhaus matches produce a \code{RiskSignal(+80)} instead of a hard
block. The signal flows through the scorer, and the dial controls whether the connection
is blocked.

\subsection{Layer 1 — ALPN Browser Bypass}
\index{bypass!ALPN browser}
\index{ALPN}
\index{h2}

Configuration key: \configkey{security\_policy.alpn\_browser\_bypass.enabled}

\begin{dangerbox}[ALPN Bypass — Do Not Disable Without Justification]
The ALPN browser bypass is structurally the most critical bypass. h2 and h1 ALPN values
are set exclusively by real browsers communicating over HTTPS. Disabling this bypass
causes all browser traffic to be scored. With typical score distributions, this will
generate false positives at scale. This bypass should only be disabled to score specific
h2 API clients, and only with monitoring in place.
\end{dangerbox}

Connections presenting \code{h2} or \code{h1} ALPN values are immediately allowed.
These connections are never enqueued for enrichment, and their IP/netblock is never
eligible for block expansion.

\subsection{Layer 1b — JA4 Whitelist}
\index{JA4!whitelist}
\index{bypass!JA4 whitelist}

Configuration key: \configkey{security\_policy.ja4\_whitelist\_bypass.enabled}

The JA4 whitelist is a Redis SET (\code{ja4:whitelist}) loaded into an in-process set
at startup and refreshed on config reload. JA4 fingerprints matching the whitelist
trigger an immediate ALLOW. Wildcard and regex matching are not supported —
whitelist entries are exact fingerprint strings.

\subsection{Layer 1c — mTLS Client Certificate}
\index{mTLS}
\index{bypass!mTLS}
\index{client certificate}

Configuration key: \configkey{security\_policy.mtls\_bypass.enabled}

A valid mTLS client certificate is cryptographic proof of identity. By default,
presenting a valid certificate against the configured trusted CA set
(\code{config/trusted\_cas.pem}) triggers an immediate ALLOW without scoring.
If disabled, mTLS clients are scored — they can still be blocked if their score
exceeds the dial threshold.

\subsection{Layer 2 — JA4 Blacklist}
\index{JA4!blacklist}
\index{bypass!JA4 blacklist}

Configuration key: \configkey{security\_policy.ja4\_blacklist\_bypass.enabled}

The JA4 blacklist is a Redis SET (\code{ja4:blacklist}) mirroring the whitelist
architecture. Matching fingerprints trigger an immediate TCP RST with no further
processing. Known blacklist entries include tool and C2 implant fingerprints such as
\fingerprint{t13d190900\_9dc949149365\_97f8aa674fd9} (Sliver C2 framework).

When disabled, blacklisted fingerprints produce a high-weight \code{RiskSignal}
routed through the scorer. This is useful for investigating traffic from
known-bad fingerprints without hard-blocking.

\section{Signal Collection (Layers 3–10)}
\index{signals!collection}
\index{asyncio.gather}

Layers 3–10 run concurrently via \code{asyncio.gather()} after the fast-path zone
completes without a decision. Each signal check is a coroutine that:
\begin{enumerate}
  \item Reads relevant data from the connection context and Redis
  \item Optionally calls an external service (AbuseIPDB, RDAP, DNS)
  \item Returns a \code{RiskSignal(score, weight, reason)} object
\end{enumerate}

All signal checks are non-blocking. External service calls are fire-and-forget via
\code{asyncio.create\_task()} — they do not block the connection handling coroutine.
Every signal check has an explicit failure handler that logs the failure, increments
a Prometheus error counter, and returns a zero/neutral result.

See Chapter~\ref{ch:signals} for full documentation of each signal.

\section{Attacker Attribution (Layer 11)}
\index{attacker attribution}
\index{JA4X}
\index{JA4T}

Layer 11 correlates JA4, JA4X, and JA4T fingerprints across IPs observed by the
analytics node. It looks for:
\begin{itemize}
  \item Multiple IPs sharing the same JA4/JA4T profile (coordinated scanning)
  \item JA4X extension ordering inconsistencies with claimed browser identity
  \item Fingerprint drift — the same IP changing fingerprints across connections
\end{itemize}

Attribution findings are written by the analytics node to \code{findings:\{ip\}} in
Redis (TTL=300s) and read back by the pipeline at layer 11.

\section{Behavioural Analysis (Layer 12)}
\index{behavioural analysis}
\index{sequential probing}
\index{coordinated bursts}

Layer 12 detects behavioural patterns that individual connections cannot reveal:
\begin{itemize}
  \item \textbf{Sequential probing} — connections to incrementally varying SNI values
        from the same source IP, characteristic of host enumeration
  \item \textbf{Coordinated bursts} — sudden increases in connection rate from multiple
        IPs sharing a fingerprint (botnet activation pattern)
  \item \textbf{Fingerprint drift} — the same IP cycling through multiple JA4 values
        to evade per-fingerprint rate limits
\end{itemize}

\section{Risk Scorer (Layer 13)}
\index{risk scorer}
\index{confidence weighting}
\index{score!0-100}

\subsection{Score Aggregation Algorithm}

The risk scorer aggregates all \code{RiskSignal} objects from layers 3–12 into a
single integer score in the range 0–100. The aggregation algorithm is
confidence-weighted:

\begin{enumerate}
  \item Each signal contributes \code{signal.score × signal.weight} to the total
  \item Weights are normalised to sum to 1.0 across all active signals
  \item The raw weighted sum is clamped to [0, 100]
  \item Signal scores from external services that are unreachable contribute 0
        (fail open — a missing signal never increases the score)
\end{enumerate}

\subsection{Score Semantics}

\begin{fullwidth}
\begin{table}[h]
\caption{Score ranges and their interpretation}
\label{tab:score-ranges}
\begin{tabular}{@{}lllp{6cm}@{}}
\toprule
\textbf{Score} & \textbf{Interpretation} & \textbf{Typical source} & \textbf{Default action at dial=100} \\
\midrule
0–19 & Clean & Browser traffic, known-good fingerprints & Allow \\
20–39 & Low risk & Datacenter IP with clean fingerprint & Allow or flag \\
40–59 & Moderate risk & Scanner, unresolved PTR, high-confidence AbuseIPDB & Rate limit \\
60–79 & High risk & Tor exit, known-bad ASN, beaconing pattern & Tarpit \\
80–100 & Critical & Confirmed C2 fingerprint, Spamhaus signal, banned IP & Block or ban \\
\bottomrule
\end{tabular}
\end{table}
\end{fullwidth}

\section{Action Decider (Layer 14)}
\index{action decider}
\index{dial}
\index{actions}

\subsection{The Dial}
\index{dial!semantics}

The \configkey{dial} value (0–100) controls how aggressively scores are translated into
blocking actions. At \code{dial=0}, every connection is allowed regardless of score —
the system is in pure monitor mode. At \code{dial=100}, the thresholds configured in
\configkey{security.rate\_limit\_strategies} apply in full.

\begin{notebox}[Dial Increment Policy]
The dial should be raised incrementally. A jump from 0 to 100 without first validating
the score distribution risks generating false positives at scale. The recommended approach
is to raise the dial to 20 and observe for 24 hours, then to 50, then to 100.
\end{notebox}

\subsection{The Six Actions}
\index{actions!allow}
\index{actions!flag}
\index{actions!rate\_limit}
\index{actions!tarpit}
\index{actions!block}
\index{actions!ban}

\begin{fullwidth}
\begin{table}[h]
\caption{Action definitions and connection handling behaviour}
\label{tab:actions}
\begin{tabular}{@{}llp{4cm}p{5.5cm}@{}}
\toprule
\textbf{Action} & \textbf{Badge} & \textbf{Connection fate} & \textbf{Notes} \\
\midrule
allow & \badgeallow & Forwarded to backend unchanged & ClientHello bytes replayed to backend; proxy becomes transparent relay \\
\addlinespace
flag & \badgeallow & Forwarded to backend & Connection allowed but high-priority event emitted to analytics stream \\
\addlinespace
rate\_limit & \badgemonitor & Queued or delayed & Token bucket or sliding window rate limiting applied per-IP or per-JA4 \\
\addlinespace
tarpit & \badgetarpit & Artificially slowed & Data delivered to backend at throttled rate; consumes attacker connection slots \\
\addlinespace
block & \badgeblock & TCP RST sent & Connection closed immediately; no data forwarded; TTL-free (not persistent) \\
\addlinespace
ban & \badgeban & TCP RST sent & IP added to Redis ban SET with TTL; all future connections blocked immediately at layer 0 \\
\bottomrule
\end{tabular}
\end{table}
\end{fullwidth}

\subsection{Ban Escalation}
\index{ban!escalation}
\index{ban!TTL}

Ban escalation applies when an IP continues to trigger blocks after an initial ban
expires. The ban TTL doubles on each re-ban up to the configured maximum:

\begin{itemize}
  \item First ban: \code{ban\_duration} seconds (default 300)
  \item Second ban: \code{ban\_duration × 2} seconds
  \item Subsequent bans: capped at 604800 seconds (7 days)
\end{itemize}

Ban state is stored in Redis as \code{ban:\{ip\}} with the TTL set at write time.
The proxy reads ban state from its local in-process cache first (TTL=30s) to avoid
Redis round-trips for known-banned IPs.

\section{Multi-Strategy Policy}
\index{multi-strategy policy}
\index{rate limiting!strategies}

Three rate-limiting strategies run simultaneously, each tracking a different key:

\begin{description}
  \item[by\_ip] Rate window keyed by client IP. Detects high-volume floods from a single address.
  \item[by\_ja4] Rate window keyed by JA4 fingerprint. Detects distributed scanning sharing a tool fingerprint.
  \item[by\_ip\_ja4\_pair] Rate window keyed by (IP, JA4) pair. Detects IP-rotation evasion where each IP presents the same fingerprint.
\end{description}

The \configkey{security.multi\_strategy\_policy} setting controls how strategy results are combined:

\begin{description}
  \item[\code{any}] Block if \emph{any} strategy triggers (most aggressive)
  \item[\code{majority}] Block if \emph{two or more} strategies trigger (default)
  \item[\code{all}] Block only if \emph{all} strategies trigger (most conservative)
\end{description}

\section{Pipeline Performance Characteristics}
\index{performance}
\index{latency}

\begin{fullwidth}
\begin{table}[h]
\caption{Expected latency contribution by pipeline zone}
\label{tab:pipeline-perf}
\begin{tabular}{@{}llll@{}}
\toprule
\textbf{Zone} & \textbf{Python (asyncio)} & \textbf{Go} & \textbf{Notes} \\
\midrule
Fast-path bypass (cache hit) & $<$100\,µs & $<$10\,µs & Local LRU cache, no Redis \\
Fast-path bypass (Redis) & 200–400\,µs & 150–300\,µs & Single Redis GET \\
Signal collection (parallel) & 1–5\,ms & 0.1–1\,ms & External services: AbuseIPDB ~50ms (cached) \\
Scoring and decision & $<$100\,µs & $<$10\,µs & Pure computation \\
Total (cached, fast path) & $<$500\,µs & $<$50\,µs & Typical cached connection \\
Total (full pipeline, no cache) & 2–10\,ms & 0.3–2\,ms & First-seen IP, all signals \\
\bottomrule
\end{tabular}
\end{table}
\end{fullwidth}

\begin{notebox}[AbuseIPDB Latency]
AbuseIPDB lookups use aggressive caching (TTL=3600s by default). A cache hit adds zero
latency. Only cache misses (first-seen IP) incur the network round-trip, which is
~50ms and runs asynchronously — it does not block the connection decision.
\end{notebox}
